<ol>
  <li>[4]
    Prove that group isomorphism is an equivalence relation, by showing that it is <ol>
      <li>Reflexive. (i.e., For all groups \(G\), \(G \cong G\).)</li>
      <li>Symmetric. (i.e., For all groups \(G_1, G_2\), if \(G_1 \cong G_2\) then \(G_2 \cong G_1\).)</li>
      <li>Transitive. (i.e., For all groups \(G_1, G_2, G_3\), if \(G_1 \cong G_2\) and \(G_2 \cong G_3\), then \(G_1 \cong G_3\).)</li>
    </ol>
  </li>
  <li>[3]
    Prove the following facts:
    <ol>
      <li>
        Every group has exactly one element of order \(1\).
      </li>
      <li>
        If \(G\) is a group of order \(p\), where \(p\) is a prime then \(G\) is a cyclic group (and thus isomorphic to the cyclic group \(C_p\)), and moreover, each non-identity element works as a generator of \(G\).
      </li>
      <li>
        If \(G\) is a group of order \(4\), then \(G\) is <i>not necessarily</i> a cyclic group, but it <i>is necessarily</i> abelian. (Hint: if \(G\) is not cyclic, then use Problem 3.1 from Exam I.)
      </li>
    </ol>
  </li>
  <li>[3]
    Here you will classify all groups up to order 11. You need to prove that all the examples you give are <i>not isomorphic</i>, but you can take for granted that I have given you the correct count.
    <br>
    It will be useful to remember that the direct product of abelian groups is abelian.
    <ol>
      <li>
        Describe the two groups of order \(4\) up to isomorphism. (You did this in the previous part!)
      </li>
      <li>
        Describe the two groups of order \(6\) up to isomorphism. (One is not abelian.)
      </li>
      <li>
        Describe the five groups of order \(8\) up to isomorphism. (Two are not abelian, and you saw one on Exam 1.)
      </li>
      <li>
        Describe the two groups of order \(9\) up to isomorphism.
      </li>
      <li>
        Describe the two groups of order \(10\) up to isomorphism. (One is not abelian.)
      </li>
      <li>
        Without looking, <i>guess</i> how many groups of order \(12\) you think there are, and how many groups of order \(16\) you think there are.
      </li>
    </ol>
  </li>
</ol>
